
\section*{\centering\LARGE{Acerca de esta publicación}}

\vspace{0.5cm}



%Partitura de este tono en formato MEI: \href{\mymeilink}{\mymeiname}

\vspace{0.5cm}

El tono aquí editado forma parte del Cancionero de Miranda. Más información sobre el cancionero en el artículo \url{http://hdl.handle.net/10261/339181}.

\vspace{0.5cm}

El resto de tonos editados se pueden encontrar en: \url{https://github.com/fernandoherreradelasheras/cancionerodemiranda/releases/}

\vspace{0.5cm}

La partitura en formato MEI (Music Encoding Initiative, un estándar abierto de codificación de música), editable en programas libres como MuseScore, está adjunta a este PDF como fichero embebido. La mayoría de lectores de PDF la muestran en su panel de adjuntos o documentos incrustados, desde donde puede guardarse para abrirla o editarla. Además, este fichero MEI, junto con los ficheros fuente de los textos y las notas, así como las herramientas para generar los pdfs finales, se encuentran en el repositorio de github \url{https://github.com/fernandoherreradelasheras/cancionerodemiranda}

\vspace{0.5cm}

BETA: Un visor interactivo de las partituras está disponible en: \url{https://cdm.humanoydivino.com/} 




